% Generated draft: requires mathematical and source review.
\documentclass[11pt]{article}
\usepackage[margin=25mm]{geometry}
\usepackage{fontspec}
\setmainfont{Latin Modern Roman}
\usepackage{amsmath,amssymb,mathtools,unicode-math}
\setmathfont{Latin Modern Math}
\usepackage{enumitem,needspace,xurl,hyperref}
\setlength{\parindent}{0pt}
\setlength{\parskip}{5pt}
\setlength{\emergencystretch}{3em}
\newcommand{\N}{\mathbb{N}}
\newcommand{\Z}{\mathbb{Z}}
\newcommand{\Q}{\mathbb{Q}}
\newcommand{\R}{\mathbb{R}}
\newcommand{\C}{\mathbb{C}}
\newcommand{\F}{\mathbb{F}}
\newcommand{\A}{\mathbb{A}}
\newcommand{\PP}{\mathbb{P}}
\newcommand{\E}{\mathbb{E}}
\newcommand{\eps}{\varepsilon}
\newcommand{\dd}{\,\mathrm{d}}
\begin{document}
\section*{UnsolvedMath ID: 30002959}
\textbf{Source problem code:} catalogue entry
\subsection{Definitions and mathematical statement}
Let \(K\) be a finite simplicial complex whose geometric realization \(|K|\) is homeomorphic to \(S^2\); such a \(K\) is a triangulated \(2\)-sphere. Under the usual Euclidean interpretation, after an ambient dimension \(d\) has been specified, a piecewise affine realization \(f\colon |K|\to\R^d\) is determined by its vertex images and is affine on each simplex. A local embedding means that every point has a neighborhood on which the restricted map is a topological embedding. The supplied record does not specify the ambient space or give the precise intended definition of a simplicial immersion.

The source asks the following questions about simplicial immersions of triangulated \(2\)-spheres:
\begin{enumerate}
\item Characterize the triangulated spheres that can be turned inside out by a discrete regular homotopy.
\item Determine the smallest triangulated sphere that admits such an eversion.
\item Determine the smallest triangulated sphere with exactly \(2^{2g}\) discrete regular homotopy classes.
\end{enumerate}

These questions retain unresolved qualifications in the supplied formulation. A precise statement requires the ambient space, the admissible immersions, the allowed deformations, whether the triangulation must remain fixed, and the endpoint condition expressing ``turned inside out.'' The first question also requires clarification of whether eversion must exist for some immersion or for every immersion of a given triangulated sphere. The meaning of ``smallest'' is not specified.

For the third question, if \(\mathcal I(K)\) denotes the admissible immersions in a specified model and \(\sim\) its specified discrete regular-homotopy equivalence relation, the requested cardinality is
\[
\#\bigl(\mathcal I(K)/{\sim}\bigr)=2^{2g}.
\]
Here \(\mathcal I(K)\) and \(\sim\) are conditional notation, not definitions supplied by the source. The parameter \(g\) has no stated meaning, domain, or quantifier. If it denotes the genus of the domain surface, then a \(2\)-sphere has \(g=0\), and the displayed cardinality becomes \(1\). The record does not establish whether that specialization is intended or whether the third question concerns surfaces of other genera. No replacement of ``sphere'' by a different surface is made in this draft.
\subsection{Short English statement}
The source asks which triangulated \(2\)-spheres admit discrete eversion, how small an evertible sphere can be, and how small a sphere with \(2^{2g}\) discrete regular homotopy classes can be. The deformation rules, size criterion, and parameter \(g\) are unspecified, so a fully determined mathematical formulation requires review.
\subsection{Sources}
\begin{itemize}
\item Local UnsolvedMath registry, supplied complete source record: permanent ID \(30002959\), title ``Minimal Triangulated Sphere Eversions,'' exact source problem code \texttt{catalogue entry}. The record belongs to the Oberwolfach Reports Open Problems collection and gives proposed year \(2015\); no individual proposer is supplied.
\item Supplied source citation: \emph{Computational Geometric and Algebraic Topology (2016)}, Oberwolfach Reports, pp.\ 2637--2699, DOI \texttt{10.4171/owr/2015/45}. \url{https://doi.org/10.4171/owr/2015/45}. The background describes Oberwolfach Reports as documenting mathematical workshops and their open-problem discussions. The supplied literature list also labels this reference ``Oberwolfach Report 45, 2015''; the differing year labels are retained without reconciliation.
\item Supplied literature reference: Kamada, \emph{Sphere Eversion from the Viewpoint of Generic Homotopy}. \url{https://doi.org/10.1016/j.topol.2017.08.005}.
\item Supplied literature reference: \emph{Oberwolfach Report 45, 2015}. \url{https://doi.org/10.4171/owr/2015/45}.
\item Literature assessment dated exactly \texttt{2026-08-21}: the supplied curation labels the problem open and states that discrete sphere eversion and the smallest triangulated sphere realizing specified discrete regular-homotopy classes remain unresolved, and that smooth sphere eversion results do not settle the discrete simplicial problem. This is supplied dated literature triage derived from the OWR curation, not an independently verified assessment; its time-sensitive status and cited claims require independent verification.
\item The registry labels the statement \texttt{corrected\_verified} and says its mathematical content was preserved while notation, typography, or omitted definitions were repaired for self-contained publication. These are supplied registry assertions, not mathematical verification performed for this draft.
\end{itemize}
\end{document}
